% GENERATED FILE. Edit canonical lessons, then run npm run generate:books.
% canonical-source-hash: fnv1a64:ecfba9dbf2f359d7
% canonical-lessons: GE-C301-kuendigen, GE-C301-kuendigung, GE-C301-befoerdern, GE-C301-befoerderung, GE-C301-wendepunkt, GE-R301-review-294-301

\chapter{Im Beruf}
\label{ch:ge-im-beruf}
\hlchaptermodality{301}

\begin{chapteropening}
\textbf{By the end of this chapter:} \emph{I can tell turning points at work: to quit, a notice, to promote, a promotion and a turning point.}

Complete the last lesson of chapter 301: I can tell turning points at work: to quit, a notice, to promote, a promotion and a turning point.
\end{chapteropening}

\section[kündigen]{kündigen — to quit, to give notice; to cancel}

\label{lesson:GE-C301-kuendigen}



\begin{quote}
Before the new one: say the German for to adopt, then the German for to go through.
\end{quote}

\subsection*{You'll want to know: kündigen}
\textbf{kündigen} — \textquotedblleft{}to quit, to give notice; to cancel\textquotedblright{}.

\begin{quote}
Sie kündigte nach zehn Jahren bei der Firma.
\end{quote}

\begin{quote}
\emph{She quit after ten years at the company.}
\end{quote}

\emph{Jemandem kündigen} — to dismiss someone. Also to cancel a \emph{Vertrag}.

\begin{cousinweb}
\emph{kündigen} is built on \emph{kund} \textquotedblleft{}known\textquotedblright{}, an inherited word with the Gothic cousin \emph{kunþs} \textquotedblleft{}known\textquotedblright{} (EtymDB): to make something known.
\end{cousinweb}

\begin{grammarlens}[title={What you've built}]
One more word for telling, explaining and describing at B1.
\end{grammarlens}

\subsection*{Your turn}
\begin{itemize}
\raggedright
  \item \emph{Say it:} \emph{kündigen}
  \item \emph{Say it:} \emph{kündigen}, once more
  \item \emph{Say it:} a sentence of your own with \emph{kündigen}
  \item \emph{Recall:} say the German for to adopt, then the German for to go through, then say \emph{kündigen} again
\end{itemize}

\begin{tcolorbox}[breakable,colback=teal!4,colframe=teal!35!black,title={Before you move on}]
What does kündigen mean? (\textquotedblleft{}To quit, to give notice; to cancel\textquotedblright{}.) Say it once more.
\end{tcolorbox}
\section[die Kündigung]{die Kündigung — the notice, the dismissal}

\label{lesson:GE-C301-kuendigung}



\begin{quote}
Before the new one: say the German for to go through, then the German for to quit.
\end{quote}

\subsection*{You'll want to know: die Kündigung}
\textbf{die Kündigung} — \textquotedblleft{}the notice, the dismissal\textquotedblright{}.

\begin{quote}
Er bekam die Kündigung am Montag.
\end{quote}

\begin{quote}
\emph{He got his notice on Monday.}
\end{quote}

Plural: \emph{die Kündigungen}.

\begin{cousinweb}
\emph{Kündigung} is built on \emph{kündigen}, from \emph{kund} \textquotedblleft{}known\textquotedblright{}, an inherited Germanic word (EtymDB).
\end{cousinweb}

\begin{grammarlens}[title={What you've built}]
One more word for telling, explaining and describing at B1.
\end{grammarlens}

\subsection*{Your turn}
\begin{itemize}
\raggedright
  \item \emph{Say it:} \emph{die Kündigung}
  \item \emph{Say it:} \emph{die Kündigung}, once more
  \item \emph{Say it:} a sentence of your own with \emph{die Kündigung}
  \item \emph{Recall:} say the German for to go through, then the German for to quit, then say \emph{die Kündigung} again
\end{itemize}

\begin{tcolorbox}[breakable,colback=teal!4,colframe=teal!35!black,title={Before you move on}]
What does die Kündigung mean? (\textquotedblleft{}The notice, the dismissal\textquotedblright{}.) Say it once more.
\end{tcolorbox}
\section[befördern]{befördern — to promote; to transport}

\label{lesson:GE-C301-befoerdern}



\begin{quote}
Before the new one: say the German for to quit, then the German for the notice.
\end{quote}

\subsection*{You'll want to know: befördern}
\textbf{befördern} — \textquotedblleft{}to promote; to transport\textquotedblright{}.

\begin{quote}
Nach zwei Jahren wurde sie befördert.
\end{quote}

\begin{quote}
\emph{After two years she was promoted.}
\end{quote}

Also 'to transport': \emph{Pakete befördern}.

\begin{cousinweb}
\emph{befördern} is built on \emph{fördern} \textquotedblleft{}to further\textquotedblright{}, from Middle High German \emph{vürdern} (EtymDB), a cousin of English \emph{further}.
\end{cousinweb}

\begin{grammarlens}[title={What you've built}]
One more word for telling, explaining and describing at B1.
\end{grammarlens}

\subsection*{Your turn}
\begin{itemize}
\raggedright
  \item \emph{Say it:} \emph{befördern}
  \item \emph{Say it:} \emph{befördern}, once more
  \item \emph{Say it:} a sentence of your own with \emph{befördern}
  \item \emph{Recall:} say the German for to quit, then the German for the notice, then say \emph{befördern} again
\end{itemize}

\begin{tcolorbox}[breakable,colback=teal!4,colframe=teal!35!black,title={Before you move on}]
What does befördern mean? (\textquotedblleft{}To promote; to transport\textquotedblright{}.) Say it once more.
\end{tcolorbox}
\section[die Beförderung]{die Beförderung — the promotion; transport}

\label{lesson:GE-C301-befoerderung}



\begin{quote}
Before the new one: say the German for the notice, then the German for to promote.
\end{quote}

\subsection*{You'll want to know: die Beförderung}
\textbf{die Beförderung} — \textquotedblleft{}the promotion; transport\textquotedblright{}.

\begin{quote}
Die Beförderung kam ganz unerwartet.
\end{quote}

\begin{quote}
\emph{The promotion came quite unexpectedly.}
\end{quote}

Plural: \emph{die Beförderungen}.

\begin{cousinweb}
\emph{Beförderung} is built on \emph{befördern}, from \emph{fördern} \textquotedblleft{}to further\textquotedblright{}, Middle High German \emph{vürdern} (EtymDB).
\end{cousinweb}

\begin{grammarlens}[title={What you've built}]
One more word for telling, explaining and describing at B1.
\end{grammarlens}

\subsection*{Your turn}
\begin{itemize}
\raggedright
  \item \emph{Say it:} \emph{die Beförderung}
  \item \emph{Say it:} \emph{die Beförderung}, once more
  \item \emph{Say it:} a sentence of your own with \emph{die Beförderung}
  \item \emph{Recall:} say the German for the notice, then the German for to promote, then say \emph{die Beförderung} again
\end{itemize}

\begin{tcolorbox}[breakable,colback=teal!4,colframe=teal!35!black,title={Before you move on}]
What does die Beförderung mean? (\textquotedblleft{}The promotion; transport\textquotedblright{}.) Say it once more.
\end{tcolorbox}
\section[der Wendepunkt]{der Wendepunkt — the turning point}

\label{lesson:GE-C301-wendepunkt}



\begin{quote}
Before the new one: say the German for to promote, then the German for the promotion.
\end{quote}

\subsection*{You'll want to know: der Wendepunkt}
\textbf{der Wendepunkt} — \textquotedblleft{}the turning point\textquotedblright{}.

\begin{quote}
Diese Reise war ein Wendepunkt in meinem Leben.
\end{quote}

\begin{quote}
\emph{This journey was a turning point in my life.}
\end{quote}

Plural: \emph{die Wendepunkte}.

\begin{cousinweb}
\emph{Wendepunkt} joins \emph{wenden} \textquotedblleft{}to turn\textquotedblright{}, Old High German \emph{wenten}, Proto-Germanic \emph{wandijaną} (EtymDB), and \emph{Punkt} \textquotedblleft{}point\textquotedblright{}, from Latin \emph{punctum}.
\end{cousinweb}

\begin{grammarlens}[title={What you've built}]
One more word for telling, explaining and describing at B1.
\end{grammarlens}

\subsection*{Your turn}
\begin{itemize}
\raggedright
  \item \emph{Say it:} \emph{der Wendepunkt}
  \item \emph{Say it:} \emph{der Wendepunkt}, once more
  \item \emph{Say it:} a sentence of your own with \emph{der Wendepunkt}
  \item \emph{Recall:} say the German for to promote, then the German for the promotion, then say \emph{der Wendepunkt} again
\end{itemize}

\begin{tcolorbox}[breakable,colback=teal!4,colframe=teal!35!black,title={Before you move on}]
What does der Wendepunkt mean? (\textquotedblleft{}The turning point\textquotedblright{}.) Say it once more.
\end{tcolorbox}
\section[(dialogue)]{Words from chapters 294–301 — a review}

\label{lesson:GE-R301-review-294-301}



\begin{quote}
Say the German for the disaster and the earthquake.
\end{quote}

\subsection*{Your turn}
Say these aloud:

\begin{itemize}
\raggedright
  \item the disaster, the earthquake, the flood, the high water, to collapse
  \item the breakdown, the breakdown service, to tow away, to damage, to destroy
  \item the birth, pregnant, the pregnancy, to baptise, the christening
  \item to grow up, starting school, the school-leaving exam, the degree, the move
  \item to get engaged, the engagement, the bride, the bridegroom, the bride and groom
  \item the marriage, the husband, the wife, the wedding ceremony, the registry office
  \item the divorce, to divorce, the separation, to adopt, to go through
  \item to quit, the notice, to promote, the promotion, the turning point
\end{itemize}

\begin{tcolorbox}[breakable,colback=teal!4,colframe=teal!35!black,title={Before you move on}]
The disaster? (\textbf{die Katastrophe}.) The earthquake? (\textbf{das Erdbeben}.) The flood? (\textbf{die Überschwemmung}.) The high water? (\textbf{das Hochwasser}.) To collapse? (\textbf{einstürzen}.) The breakdown? (\textbf{die Panne}.) The breakdown service? (\textbf{der Pannendienst}.) To tow away? (\textbf{abschleppen}.) To damage? (\textbf{beschädigen}.) To destroy? (\textbf{zerstören}.) The birth? (\textbf{die Geburt}.) Pregnant? (\textbf{schwanger}.) The pregnancy? (\textbf{die Schwangerschaft}.) To baptise? (\textbf{taufen}.) The christening? (\textbf{die Taufe}.) To grow up? (\textbf{aufwachsen}.) Starting school? (\textbf{die Einschulung}.) The school-leaving exam? (\textbf{das Abitur}.) The degree? (\textbf{der Abschluss}.) The move? (\textbf{der Umzug}.) To get engaged? (\textbf{verloben}.) The engagement? (\textbf{die Verlobung}.) The bride? (\textbf{die Braut}.) The bridegroom? (\textbf{der Bräutigam}.) The bride and groom? (\textbf{das Brautpaar}.) The marriage? (\textbf{die Ehe}.) The husband? (\textbf{der Ehemann}.) The wife? (\textbf{die Ehefrau}.) The wedding ceremony? (\textbf{die Trauung}.) The registry office? (\textbf{das Standesamt}.) The divorce? (\textbf{die Scheidung}.) To divorce? (\textbf{scheiden}.) The separation? (\textbf{die Trennung}.) To adopt? (\textbf{adoptieren}.) To go through? (\textbf{erleben}.) To quit? (\textbf{kündigen}.) The notice? (\textbf{die Kündigung}.) To promote? (\textbf{befördern}.) The promotion? (\textbf{die Beförderung}.) The turning point? (\textbf{der Wendepunkt}.)
\end{tcolorbox}
